% Generated draft: requires mathematical and source review.
\documentclass[11pt]{article}
\usepackage[margin=25mm]{geometry}
\usepackage{fontspec}
\setmainfont{Latin Modern Roman}
\usepackage{amsmath,amssymb,mathtools,unicode-math}
\setmathfont{Latin Modern Math}
\usepackage{enumitem,needspace,xurl,hyperref}
\setlength{\parindent}{0pt}
\setlength{\parskip}{5pt}
\setlength{\emergencystretch}{3em}
\newcommand{\N}{\mathbb{N}}
\newcommand{\Z}{\mathbb{Z}}
\newcommand{\Q}{\mathbb{Q}}
\newcommand{\R}{\mathbb{R}}
\newcommand{\C}{\mathbb{C}}
\newcommand{\F}{\mathbb{F}}
\newcommand{\A}{\mathbb{A}}
\newcommand{\PP}{\mathbb{P}}
\newcommand{\E}{\mathbb{E}}
\newcommand{\eps}{\varepsilon}
\newcommand{\dd}{\,\mathrm{d}}
\begin{document}
\section*{UnsolvedMath ID: 20000394}
\textbf{Source problem code:} catalogue entry
\subsection{Definitions and mathematical statement}
\textbf{Original question (Brian Conrey).} Fix a prime power $q$. As $g\to\infty$, describe the distribution of discriminants of $L$-polynomials of curves, or of hyperelliptic curves, over $\F_q$. The sampling measure, smoothness/projectivity conventions, signed versus absolute discriminant, and normalization are not specified.

\textbf{Attached unreviewed reduction.} In the smooth projective geometrically connected interpretation, write
\[Z_C(T)=\frac{L_C(T)}{(1-T)(1-qT)},\qquad L_C(T)=\prod_{i=1}^{2g}(1-\alpha_iT),\qquad P_C(T)=T^{2g}L_C(T^{-1})=\prod_i(T-\alpha_i),\]
with $|\alpha_i|=\sqrt q$. The aid asserts equality of the discriminants of the nonmonic $L_C$ and monic reciprocal $P_C$, using the usual polynomial-discriminant convention including leading coefficients. Put $N=2g$, $z_i=\alpha_i/\sqrt q$, and
\[D_C=q^{-g(2g-1)}\operatorname{Disc}(P_C)=\prod_{i<j}(z_i-z_j)^2.\]
For a real Weil polynomial $h_C$ of degree $g$, defined by $P_C(T)=T^gh_C(T+q/T)$ with roots $\beta_j\in[-2\sqrt q,2\sqrt q]$, the reported factorization is
\[\operatorname{Disc}(P_C)=q^{g(g-1)}\operatorname{Disc}(h_C)^2\prod_{j=1}^g(\beta_j^2-4q).\]
It records a zero precisely when $h_C$ has a repeated root or some $\beta_j=\pm2\sqrt q$; off this locus the sign is $(-1)^g$. The square class in $\Q^\times/\Q^{\times2}$ is that of $h_C(2\sqrt q)h_C(-2\sqrt q)$, since the other factors are squares. These formulas require the standard Weil-polynomial conventions and independent review.

On the nonzero locus let $t_m=\sum_i z_i^m$, the real normalized Frobenius trace. The aid's Abel trace-square reduction reads
\[\log|D_C|=\lim_{r\uparrow1}\sum_{m\ge1}\frac{r^m}{m}(N-t_m^2).\]
It reports an explicit finite truncation estimate controlled by the minimum eigenvalue gap $\min_{i\ne j}|z_i-z_j|$. The actual bound and truncation parameters are not supplied, so no uniform small-gap control is inferred. The repeated-root locus contributes an atom at discriminant zero and cannot be handled by this logarithmic formula.

The aid also reports a sharp support bound for the normalized absolute discriminant and an exact Haar-$\mathrm{USp}(2g)$ Selberg moment law, but supplies neither formula or reference in the summary. The Haar law is explicitly only a benchmark: it is not identified with the fixed-$q$ arithmetic distribution. The proposed reduction leaves the repeated-root mass, minimum-gap tails and growing-range joint squared-trace statistics as the unresolved arithmetic inputs. It does not prove a limiting law for curves or hyperelliptic curves.
\subsection{Short English statement}
The problem asks for a fixed-field large-genus discriminant law. The aid reduces normalized discriminants to a real Weil polynomial and, away from repeated roots, to an Abel sum of squared traces. Zero mass and small eigenvalue gaps remain essential inputs, while Haar symplectic moments are retained only as a comparison model.
\subsection{Sources}
\begin{itemize}
\item Local UnsolvedMath registry, permanent ID 20000394, source code catalogue entry; \texttt{aim-algebraic-number-theory-notes.json}, notes[45], tag \texttt{question}.
\item AIM workshop \emph{Arithmetic statistics over finite fields and function fields}, Question 16 (Brian Conrey): \url{https://aimath.org/pastworkshops/arithstatffieldproblems.pdf}. No original remarks or bibliography are supplied; no PDF was consulted.
\item Attached machine-generated research attempt 1: result \texttt{reduction}, problem status at run \texttt{open}, rigor label \texttt{proved}, novelty confidence low. It is not peer reviewed and requires independent verification before citation. No authoritative citation is supplied for the reductions, support bound or Haar moment law.
\end{itemize}
\end{document}
